\documentclass[12pt, a4paper]{article}

% ── Packages ──────────────────────────────────────────────────────────────────
\usepackage[utf8]{inputenc}
\usepackage[T1]{fontenc}
\usepackage{lmodern}
\usepackage{amsmath, amssymb, amsthm}
\usepackage{mathtools}          % for \coloneqq, \DeclarePairedDelimiter, etc.
\usepackage{booktabs}           % nicer tables
\usepackage{enumitem}           % custom list formatting
\usepackage{hyperref}           % clickable cross-references
\usepackage[margin=2.5cm]{geometry}
\usepackage{parskip}            % blank lines between paragraphs, no indent
\usepackage{xcolor}
\usepackage{tcolorbox}          % coloured boxes for "intuitive explanation" etc.
\usepackage{titlesec}           % section formatting

% ── Colours ───────────────────────────────────────────────────────────────────
\definecolor{intuitionbg}{RGB}{240,248,255}   % pale blue
\definecolor{simulationbg}{RGB}{255,250,240}  % pale orange
\definecolor{keybox}{RGB}{245,245,220}        % pale yellow

% ── Environments ──────────────────────────────────────────────────────────────
\newtcolorbox{intuition}{
  colback=intuitionbg, colframe=blue!40!black,
  title={\textbf{Intuitive Explanation}}, fonttitle=\bfseries
}
\newtcolorbox{simulation}{
  colback=simulationbg, colframe=orange!60!black,
  title={\textbf{Simulation}}, fonttitle=\bfseries
}
\newtcolorbox{keytakeaway}{
  colback=keybox, colframe=olive!60!black,
  title={\textbf{Key Takeaway}}, fonttitle=\bfseries
}
\newtcolorbox{plainlanguage}{
  colback=white, colframe=gray!50,
  left=6pt, right=6pt, top=4pt, bottom=4pt
}

\newtheorem{theorem}{Theorem}
\newtheorem{definition}{Definition}

% ── Notation shorthand ────────────────────────────────────────────────────────
\newcommand{\E}{\mathbb{E}}            % expected value
\newcommand{\Prob}{\mathbb{P}}         % probability
\newcommand{\Var}{\operatorname{Var}}
\newcommand{\Cov}{\operatorname{Cov}}
\newcommand{\SD}{\operatorname{SD}}
\newcommand{\given}{\,\vert\,}         % conditional bar with spacing
\newcommand{\LR}{\Lambda}              % likelihood ratio


% ══════════════════════════════════════════════════════════════════════════════
\title{\textbf{Presentation Plan: From Random Walks to Sequential Testing}\\[6pt]
       \large Version~3 — \LaTeX{} Edition}
\author{}
\date{March 2026}

\begin{document}
\maketitle
\tableofcontents
\newpage

% ══════════════════════════════════════════════════════════════════════════════
\section*{About This Document}
\addcontentsline{toc}{section}{About This Document}

\textbf{Target audience.}
People comfortable with algebra and basic probability (e.g.\ a 12th-grade
math background).  No calculus is required to follow the main ideas, though
some calculus notation (integrals) appears in later acts for completeness ---
always with plain-language translations.

\textbf{Format.}
Eight interactive acts --- deliberately small steps.  Each act has a
simulation, an intuitive explanation, and a mathematical section.  The math
sections always define every symbol before using it.  Every formula is
accompanied by a plain-language translation immediately below it.

\textbf{Design philosophy.}
Better to over-explain than under-explain.  When in doubt, slow down.

% ══════════════════════════════════════════════════════════════════════════════
\section{Notation Cheat Sheet}
\label{sec:notation}

This cheat sheet should appear as a collapsible sidebar in the actual
webpage, accessible from every act.

\begin{center}
\begin{tabular}{@{} l l p{8.5cm} @{}}
\toprule
\textbf{Symbol}       & \textbf{Name}                  & \textbf{What it means} \\
\midrule
$X,\,M,\,Y$          & Capital letters                & Quantities that are random \\
$X_1, X_2, X_n$      & Subscript numbers              & ``The value of $X$ at step 1, 2, $n$'' \\
$\E[\cdot]$           & Expected value                 & Long-run average if you repeated the process many times \\
$\E[\cdot \given \cdot]$ & Conditional expectation     & Long-run average, \emph{given that you already know} whatever appears after the $\vert$ \\
$\Prob(\cdot)$        & Probability                    & A number between 0 and 1 \\
$\displaystyle\sum$   & Summation ($\Sigma$)           & ``Add all of these up'' \\
$\displaystyle\prod$  & Product ($\Pi$)                & ``Multiply all of these together'' \\
$\displaystyle\int$   & Integral                       & Continuous version of summation --- a smooth $\sum$ \\
$\geq,\;\leq$        & Greater/less than or equal to   & \\
$\infty$              & Infinity                       & A process that goes on forever \\
$\sqrt{\cdot}$        & Square root                    & $\sqrt{9}=3$, $\sqrt{100}=10$ \\
$\exp(x)$            & Exponential function            & $e\approx 2.718$ raised to the power $x$ \\
$\sup$                & Supremum                       & The highest value something ever reaches \\
$\alpha$              & Alpha                          & False positive rate \\
$\beta$               & Beta                           & False negative rate \\
$\delta$              & Delta                          & Effect size \\
$\sigma$              & Sigma (lowercase)              & Standard deviation of individual observations \\
$\tau$                & Tau                            & In mSPRT: width of the prior over effect sizes \\
$\theta$              & Theta                          & In CUPED: variance reduction coefficient \\
$\rho$                & Rho                            & Correlation coefficient ($-1$ to $+1$) \\
$\LR$                 & Lambda (uppercase)             & Likelihood ratio --- a running score of evidence \\
\bottomrule
\end{tabular}
\end{center}


% ══════════════════════════════════════════════════════════════════════════════
% ACT 0
% ══════════════════════════════════════════════════════════════════════════════
\newpage
\section{Act~0 --- Why Should You Care?  The Peeking Problem}

\subsection{The Motivation}

\begin{intuition}
You work at a company running an A/B test.  Half your users see the old
website, half see a new design.  Every day, you check the dashboard:
``Is the new design better?''  After a week, the dashboard shows $p=0.03$ ---
that looks significant!  You call it: the new design wins.

But here is the dirty secret of traditional statistics: \textbf{the more often
you check, the more likely you are to see a false positive.}  A test designed
to have a 5\% false positive rate can easily reach 20--30\% if you peek at it
repeatedly.  One in four or five ``winning'' experiments might actually be noise.

This is not a theoretical concern.  It is the single most common statistical
mistake in online experimentation.
\end{intuition}

\subsection{The Simulation}

\begin{simulation}
Two groups of observations (control and treatment) are generated from
\emph{identical} distributions --- there is no real effect.  A standard
$t$-test is computed after every new observation arrives.

A counter shows: ``Number of times $p<0.05$ so far.''  The user watches the
$p$-value bounce around.  Every time it dips below $0.05$, it lights up red ---
a false positive.

\textbf{Users can control:}
\begin{itemize}[nosep]
  \item How many total observations to simulate
  \item How many independent experiments to run in parallel
\end{itemize}

A bar chart accumulates: ``Fraction of experiments that showed $p<0.05$ at
\emph{some} point during peeking.''  For 1,000 simulated experiments with
regular peeking, this fraction climbs well above 5\% --- often to 20--30\%.

A second mode checks only once, at the very end.  That fraction stays near 5\%.
\end{simulation}

\subsection{Intuitive Explanation}

\begin{intuition}
Traditional statistical tests promise: ``If there is no real effect, I will
falsely cry wolf at most 5\% of the time.''  But this promise comes with fine
print: \textbf{you may only look at the result once, at a pre-determined time.}

Every time you peek, you get another chance for randomness to fool you.  It is
like rolling a die repeatedly --- the more rolls, the more likely you are to
get a six eventually, even though any single roll has only a $1/6$ chance.

The rest of this presentation builds a test that does not have this problem:
a test where you can peek as often as you like and the false positive guarantee
still holds.

That promise has a name: \textbf{anytime-valid inference}.  Getting there
requires some beautiful mathematics, which we will build up one step at a
time.
\end{intuition}

\begin{keytakeaway}
\textbf{The peeking problem:} Checking a traditional test repeatedly inflates
the false positive rate far beyond its nominal level.  We need a fundamentally
different kind of test --- one that remains valid no matter when or how often
we look.
\end{keytakeaway}


% ══════════════════════════════════════════════════════════════════════════════
% ACT 1
% ══════════════════════════════════════════════════════════════════════════════
\newpage
\section{Act~1 --- The Random Walk}

\subsection{The Simulation}

\begin{simulation}
A vertical number line appears on screen.  A dot starts at position~0.  Each
step, a coin is flipped.  Heads: the dot moves up one step.  Tails: down one
step.  The path is drawn as a graph (position vs.\ step number).

\textbf{User controls:} speed, number of steps, probability of going up
(initially fixed at $50/50$; later a slider), number of simultaneous paths.
\end{simulation}

\subsection{Intuitive Explanation}

\begin{intuition}
Imagine you are standing in the middle of a long hallway.  Flip a coin.
Heads --- step forward.  Tails --- step back.  Repeat hundreds of times.

Where will you end up?  We don't know exactly, but we can say something
precise about the \emph{range} of likely outcomes.  Most of the time you won't
wander far.  Occasionally, a long streak of heads or tails carries you far in
one direction.  The longer you walk, the wider the range of possible positions
--- but the centre of that range stays at zero.

When we run many people simultaneously, they fan out over time --- like ink
dropped into water.  This spreading has a precise mathematical shape.
\end{intuition}

\subsection{Mathematical Formulation}

\subsubsection*{Setting up the notation}

Let $X_i$ represent the outcome of the $i$th coin flip:
\[
  X_i =
  \begin{cases}
    +1 & \text{if heads (step forward),} \\
    -1 & \text{if tails (step backward).}
  \end{cases}
\]
The probability of heads is written $p$.  For a fair coin, $p = 0.5$.

\subsubsection*{Position after $n$ steps}

\[
  \boxed{S_n = \sum_{i=1}^{n} X_i = X_1 + X_2 + \cdots + X_n}
\]

\begin{plainlanguage}
Read as: ``$S$ sub $n$ equals the sum of all $X_i$, starting from $i=1$ up to
$i=n$.''
\end{plainlanguage}

\textbf{Worked example.}  First 5 flips: H, T, H, H, T, so
$X_1\!=\!+1,\; X_2\!=\!-1,\; X_3\!=\!+1,\; X_4\!=\!+1,\; X_5\!=\!-1$.
\[
  S_1=1,\quad S_2=0,\quad S_3=1,\quad S_4=2,\quad S_5=1.
\]

\subsubsection*{The average outcome (expected value)}

When $p = 0.5$, each flip has expected value zero:
\[
  \E[X_i]
    = \underbrace{0.5}_{P(\text{heads})} \times (+1)
    + \underbrace{0.5}_{P(\text{tails})} \times (-1)
    = 0.
\]
Because each step averages to zero:
\[
  \boxed{\E[S_n] = \E[X_1] + \E[X_2] + \cdots + \E[X_n] = 0}
\]

\begin{plainlanguage}
On average, across many trials, you end up back where you started.
\end{plainlanguage}

\subsubsection*{How spread out the outcomes are --- variance \& standard deviation}

\textbf{Variance} measures the average squared distance from the mean.
For a single fair coin flip:
\[
  \Var(X_i) = \E\bigl[(X_i - 0)^2\bigr]
  = 0.5 \times (+1)^2 + 0.5 \times (-1)^2
  = 1.
\]
Because flips are independent, the variance of the sum is the sum of
the variances:
\[
  \Var(S_n) = \Var(X_1) + \Var(X_2) + \cdots + \Var(X_n)
  = \underbrace{1 + 1 + \cdots + 1}_{n\text{ terms}} = n.
\]
The \textbf{standard deviation} is the square root of the variance ---
same units as the position:
\[
  \boxed{\SD(S_n) = \sqrt{\Var(S_n)} = \sqrt{n}}
\]

\begin{plainlanguage}
After 100 steps, the typical distance from zero is $\sqrt{100}=10$.
After 10,000 steps, it is $\sqrt{10{,}000}=100$.  The spread grows with the
\emph{square root} of the number of steps, not the number itself.
\end{plainlanguage}

The simulation shows the $\pm\sqrt{n}$ envelope as two curved lines.

\subsubsection*{What if the coin is not fair?}

When $p \neq 0.5$, each step has a non-zero average:
\[
  \E[X_i] = p \times (+1) + (1-p) \times (-1) = 2p - 1.
\]
For $p = 0.6$: $\E[X_i] = 0.2$.  After $n$ steps:
\[
  \E[S_n] = n(2p-1).
\]
This is a systematic \emph{drift} --- the random walk now has a trend.

\begin{keytakeaway}
\textbf{Key concepts:} stochastic process, expected value, variance, standard
deviation, the $\sqrt{n}$ growth of randomness.
\end{keytakeaway}


% ══════════════════════════════════════════════════════════════════════════════
% ACT 2
% ══════════════════════════════════════════════════════════════════════════════
\newpage
\section{Act~2 --- The Gambling Game and the Martingale}

\subsection{The Simulation}

\begin{simulation}
The same random walk, relabelled.  The vertical axis now shows
``\texteuro{} profit or loss.''  A bold horizontal line marks zero ---
break-even.

\textbf{Probability slider:} $p > 0.5$ gives the player an edge;
$p < 0.5$ gives the house an edge.

\textbf{Doubling strategy mode:} Runs 10,000~gamblers using the martingale
betting strategy.  A histogram shows the distribution of final wealth --- many
small winners, a few catastrophic losers.  A counter shows average profit
converging to zero.

\textbf{Build your own stopping rule:} Users define a rule (``Stop when ahead
by \texteuro$K$'').  The simulation runs 10,000~gamblers.  Average winnings at
stopping converge to zero.
\end{simulation}

\subsection{Intuitive Explanation}

\begin{intuition}
Win \texteuro1 for heads, lose \texteuro1 for tails.  Cumulative winnings
trace the same random walk from Act~1.

The gambler's temptation: ``I'll keep playing until I'm ahead, then quit.''
Sounds foolproof --- surely a run of luck will eventually put you in the green?

Mathematically, a fair random walk does eventually return to positive territory.
But the time it takes can be astronomically long, and while you wait you might
lose everything.

\textbf{Central insight:} There is no strategy that improves your expected
outcome in a fair game.  This is a proven mathematical theorem.  Try any
stopping rule in the simulation --- the average winnings always converge to
zero.
\end{intuition}

\begin{intuition}
\textbf{The doubling strategy (Martingale betting), exposed.}

Lose, double the bet, repeat until you win back your losses plus \texteuro1.
It works --- most of the time.

But a long losing streak forces bets of
\texteuro1, 2, 4, 8, 16, 32, 64, \ldots\
After just 10 consecutive losses, you must bet \texteuro1,024 to win back
\texteuro1.  The histogram makes it vivid: a tall bar of people who won
\texteuro1, and a barely visible bar at the far left of people who lost
thousands.  Multiply heights by positions and sum: the average is zero.
\end{intuition}

\subsection{Mathematical Formulation}

\begin{definition}[Martingale]
A sequence $M_0, M_1, M_2, \ldots$ is a \textbf{martingale} if, at every
step, the best prediction of the next value is the current value:
\[
  \boxed{\E\bigl[M_n \given M_0, M_1, \ldots, M_{n-1}\bigr] = M_{n-1}}
\]
\end{definition}

\textbf{Reading the notation carefully:}
\begin{itemize}[nosep]
  \item $\E[\cdot]$ = ``the expected value of\ldots'' (long-run average).
  \item The vertical bar ``$\given$'' = ``given that we know.''
        Everything after the bar is information we have.
  \item $M_0, M_1, \ldots, M_{n-1}$ = the full history up to step $n-1$.
  \item Full sentence: ``The expected value of $M_n$, given that we know
        everything up to step $n-1$, equals $M_{n-1}$.''
\end{itemize}

\begin{plainlanguage}
Knowing the entire history, your best guess for the next value is just the
current value.  The past gives no useful information about the future direction.
\end{plainlanguage}

\subsubsection*{Verifying: the coin-flip game is a martingale}

Let $M_n = S_n$ (cumulative winnings).  Then $M_n = M_{n-1} + X_n$.
\begin{align*}
  \E[M_n \given M_0,\ldots,M_{n-1}]
  &= \E[M_{n-1} + X_n \given M_0,\ldots,M_{n-1}] \\
  &= M_{n-1} + \E[X_n \given M_0,\ldots,M_{n-1}]
     & &\text{($M_{n-1}$ is known, passes through $\E$)} \\
  &= M_{n-1} + 0
     & &\text{(next flip is independent; $\E[X_n]=0$)} \\
  &= M_{n-1}. \qquad\checkmark
\end{align*}

\subsubsection*{Doob's Optional Stopping Theorem (1953)}

\begin{theorem}[Doob's Optional Stopping Theorem]
Let $\{M_n\}$ be a martingale and $\tau$ a \textbf{stopping time} --- a rule
for when to quit that uses only information accumulated so far (no looking into
the future).  Then, under reasonable conditions:
\[
  \boxed{\E[M_\tau] = \E[M_0]}
\]
\end{theorem}

\begin{plainlanguage}
For a gambler starting with $M_0 = 0$: expected winnings at stopping $= 0$,
always, for any strategy.  \textbf{You cannot beat a fair game by choosing
when to quit.}
\end{plainlanguage}

\subsubsection*{What is a stopping time?}

$\tau$ is a decision rule of the form: ``Stop after step $n$ if [condition
based only on what happened so far].''

\begin{center}
\begin{tabular}{@{} l c @{}}
\toprule
\textbf{Rule} & \textbf{Valid?} \\
\midrule
Stop after exactly 100 flips & \checkmark \\
Stop the first time I am ahead by \texteuro10 & \checkmark \\
Stop when I have lost \texteuro50 & \checkmark \\
Stop after 1000 flips \emph{or} when ahead by \texteuro5, whichever first & \checkmark \\
Stop right before the next loss & $\times$ (requires the future!) \\
\bottomrule
\end{tabular}
\end{center}

\subsubsection*{The ``reasonable conditions'' for Doob's theorem}

Doob's theorem requires at least one of:
\begin{enumerate}[nosep]
  \item $\tau \leq N$ for some fixed $N$ (bounded stopping time), \textbf{or}
  \item $|M_n| \leq C$ for some fixed $C$ (bounded martingale), \textbf{or}
  \item $\E[\tau] < \infty$ (finite expected stopping time).
\end{enumerate}

When \emph{none} hold --- as with the doubling strategy --- the guarantee
breaks down.

\subsubsection*{Why the doubling strategy fails}

The doubling strategy guarantees eventual profit \emph{only if} you have
infinite money and infinite time.  In practice:
\begin{itemize}[nosep]
  \item Finite budget $\Rightarrow$ real probability of ruin before recovery.
  \item Expected time to recovery is infinite: $\E[\tau] = \infty$.
  \item Doob's theorem gives $\E[M_\tau] \leq 0$, not $= 0$.
\end{itemize}

\subsubsection*{Why this matters for what is coming}

The martingale concept reappears in every subsequent act.  The likelihood
ratio (Act~4) is a martingale.  The mSPRT statistic (Act~7) is a martingale.
The reason sequential tests control false positives --- even with peeking ---
is that they exploit the same structure: \textbf{a fair game cannot be beaten
by choosing when to stop.}

\begin{keytakeaway}
\textbf{Key concepts:} martingale, conditional expectation, stopping time,
Doob's Optional Stopping Theorem, the ``reasonable conditions.''
\end{keytakeaway}


% ══════════════════════════════════════════════════════════════════════════════
% ACT 3
% ══════════════════════════════════════════════════════════════════════════════
\newpage
\section{Act~3 --- Hypothesis Testing: Telling Two Coins Apart}

\subsection{The Simulation}

\begin{simulation}
The user is presented with a coin and must decide: fair or biased?

Two panels run side by side:
\begin{itemize}[nosep]
  \item \textbf{Left:} Coin is actually fair ($p=0.5$).
  \item \textbf{Right:} Coin is actually biased ($p=0.5+\delta$).
\end{itemize}
After each flip, a running tally shows heads/tails and current proportion.
For small $\delta$, the sequences look nearly identical for many flips.

\textbf{User controls:} true bias $\delta$; number of flips before deciding.
\end{simulation}

\subsection{Intuitive Explanation}

\begin{intuition}
You suspect a coin might be unfair.  Flip it many times and look at the
results.  10~flips with 6~heads?  Could be luck.  1,000~flips with 600~heads?
Strong evidence.

Statisticians frame this as two competing stories:
\begin{itemize}[nosep]
  \item \textbf{Null hypothesis $H_0$:} ``The coin is fair.  Nothing is happening.''
  \item \textbf{Alternative hypothesis $H_1$:} ``The coin is biased.  Something real is going on.''
\end{itemize}

Two types of mistakes:
\begin{enumerate}[nosep]
  \item \textbf{False positive} (Type~I): conclude biased when it is fair.
  \item \textbf{False negative} (Type~II): conclude fair when it is biased.
\end{enumerate}
Traditionally: $\alpha = 0.05$ (max false positive rate),
$\beta = 0.20$ (max false negative rate).
\textbf{Power} $= 1 - \beta = 0.80$.
\end{intuition}

\subsection{Mathematical Formulation}

\subsubsection*{The likelihood}

Suppose $n$ flips produce $k$ heads.

Under $H_0$ (fair coin, $p = 0.5$):
\[
  \mathcal{L}_0
  = (0.5)^k \times (0.5)^{n-k}
  = (0.5)^n.
\]

Under $H_1$ (biased coin, $p = 0.5 + \delta$):
\[
  \mathcal{L}_1
  = (0.5 + \delta)^k \times (0.5 - \delta)^{n-k}.
\]

\subsubsection*{The likelihood ratio}

\[
  \boxed{\LR_n
  = \frac{\mathcal{L}_1}{\mathcal{L}_0}
  = \frac{(0.5+\delta)^k \;(0.5-\delta)^{n-k}}{(0.5)^n}}
\]

\begin{plainlanguage}
``How many times more likely is our observed data under the biased-coin story
compared to the fair-coin story?''
\end{plainlanguage}

\textbf{Worked example.}
$\delta = 0.1$ (biased coin has $p=0.6$); $n=10$ flips, $k=7$ heads.
\begin{align*}
  \mathcal{L}_1 &= (0.6)^7 \times (0.4)^3 = 0.0280 \times 0.064 \approx 0.001792, \\
  \mathcal{L}_0 &= (0.5)^{10} \approx 0.000977, \\
  \LR_{10} &= \frac{0.001792}{0.000977} \approx 1.83.
\end{align*}
The data is $\approx\!1.83$ times more likely if the coin is biased.
Evidence, but not overwhelming.

\subsubsection*{Interpreting $\LR_n$}

\begin{center}
\begin{tabular}{@{} r l @{}}
\toprule
$\LR_n = 1$   & Data equally consistent with both hypotheses \\
$\LR_n = 10$  & 10$\times$ more likely under $H_1$ --- moderate evidence \\
$\LR_n = 100$ & 100$\times$ more likely under $H_1$ --- strong evidence \\
$\LR_n = 0.1$ & 10$\times$ more likely under $H_0$ --- evidence for fairness \\
\bottomrule
\end{tabular}
\end{center}

\subsubsection*{Incremental updating}

$\LR_n$ can be computed one flip at a time:
\[
  \boxed{\LR_n = \LR_{n-1} \times \frac{f_1(x_n)}{f_0(x_n)}}
\]
where $f_j(x_n)$ is the probability of the $n$th flip's outcome under $H_j$.
\begin{itemize}[nosep]
  \item Heads: multiply by $\dfrac{0.5+\delta}{0.5}$.
  \item Tails: multiply by $\dfrac{0.5-\delta}{0.5}$.
\end{itemize}
This incremental structure is what makes \emph{sequential} testing possible.

\subsubsection*{Starting value}

Before any data:
\[
  \LR_0 = 1.
\]
Neither hypothesis is favoured.  This will be important for Ville's inequality
(which requires the starting value to equal~1).

\begin{keytakeaway}
\textbf{Key concepts:} null hypothesis, alternative hypothesis, likelihood,
likelihood ratio, false positive, false negative, power, incremental updating.
\end{keytakeaway}


% ══════════════════════════════════════════════════════════════════════════════
% ACT 4
% ══════════════════════════════════════════════════════════════════════════════
\newpage
\section{Act~4 --- The Likelihood Ratio Is a Martingale}

\subsection{The Simulation}

\begin{simulation}
The likelihood ratio $\LR_n$ is now plotted as a path over time --- like the
random walk from Act~1.

\textbf{Left panel:} 100 paths of $\LR_n$ when the coin is truly \emph{fair}
($H_0$ true).  Paths wander with no systematic trend.

\textbf{Right panel:} 100 paths when the coin is truly \emph{biased}
($H_1$ true).  Paths drift upward.
\end{simulation}

\subsection{Intuitive Explanation}

\begin{intuition}
Under $H_0$, the likelihood ratio wanders up and down but never develops a
consistent trend.  On average, it stays at~1.

This \emph{is} a martingale --- exactly like the gambler's cumulative winnings.
Everything from Act~2 applies: Doob's theorem says you cannot systematically
make $\LR_n$ large by choosing clever peeking times.
\end{intuition}

\subsection{Mathematical Formulation}

\subsubsection*{Proof that $\LR_n$ is a martingale under $H_0$}

We must show:
$\E[\LR_n \given \LR_0, \ldots, \LR_{n-1}] = \LR_{n-1}$.

From Act~3, $\LR_n = \LR_{n-1} \times \dfrac{f_1(x_n)}{f_0(x_n)}$.
Therefore:
\begin{align}
  \E\bigl[\LR_n \given \text{past}\bigr]
  &= \E\!\left[\LR_{n-1} \cdot \frac{f_1(x_n)}{f_0(x_n)} \;\bigg\vert\; \text{past}\right]
     \nonumber \\
  &= \LR_{n-1} \cdot \E\!\left[\frac{f_1(x_n)}{f_0(x_n)}\right]
     & &\text{($\LR_{n-1}$ is known)}
     \label{eq:lr-mart-step1}
\end{align}

Now the key step.  Under $H_0$, $x_n$ is drawn from $f_0$.  So:
\begin{align}
  \E\!\left[\frac{f_1(x_n)}{f_0(x_n)}\right]
  &= \sum_{\text{all outcomes } x} f_0(x) \cdot \frac{f_1(x)}{f_0(x)}
     \nonumber \\[6pt]
  &= \sum_x \cancel{f_0(x)} \cdot \frac{f_1(x)}{\cancel{f_0(x)}}
     & &\text{\textbf{(the cancellation!)}}
     \nonumber \\[6pt]
  &= \sum_x f_1(x)
     \nonumber \\[4pt]
  &= 1.
     & &\text{($f_1$ is a probability distribution; probabilities sum to~1)}
     \label{eq:lr-mart-cancel}
\end{align}

Therefore:
\[
  \E\bigl[\LR_n \given \text{past}\bigr]
  = \LR_{n-1} \times 1
  = \LR_{n-1}. \qquad\checkmark
\]

\subsubsection*{The cancellation with actual numbers}

For $\delta = 0.1$ (biased coin $p=0.6$):

\medskip
\begin{center}
\begin{tabular}{@{} l c c c @{}}
\toprule
\textbf{Outcome} & $f_0(x)$ & $f_1(x)/f_0(x)$ & $f_0(x)\times\bigl[f_1(x)/f_0(x)\bigr]$ \\
\midrule
Heads & 0.5 & $0.6/0.5 = 1.2$ & $0.5 \times 1.2 = 0.6$ \\
Tails & 0.5 & $0.4/0.5 = 0.8$ & $0.5 \times 0.8 = 0.4$ \\
\midrule
\multicolumn{3}{r}{\textbf{Total:}} & $\mathbf{1.0}$ \checkmark \\
\bottomrule
\end{tabular}
\end{center}

\begin{plainlanguage}
When we compute the expected ratio under $H_0$, the $H_0$ probabilities
cancel out, and we are just left summing $H_1$ probabilities --- which must
add to~1.  This is the ``aha'' moment of the proof.
\end{plainlanguage}

\subsubsection*{Two additional properties}

\begin{enumerate}[nosep]
  \item \textbf{$\LR_n$ is non-negative:} It is a product of non-negative
        ratios (probabilities divided by probabilities), so $\LR_n \geq 0$.
  \item \textbf{$\LR_0 = 1$:} Before any data, neither hypothesis is favoured.
\end{enumerate}

So $\LR_n$ is a \textbf{non-negative martingale starting at~1}.  Remember
these properties --- they are exactly what Ville's inequality needs.

\subsubsection*{Under $H_1$, $\LR_n$ drifts upward}

When the coin really is biased:
\[
  \E\!\left[\frac{f_1(x_n)}{f_0(x_n)}\right]_{\!H_1}
  = \sum_x f_1(x) \cdot \frac{f_1(x)}{f_0(x)}
  > 1.
\]
So $\LR_n$ tends to grow --- this upward drift is the signal the test detects.

\begin{keytakeaway}
\textbf{Key concepts:} the likelihood ratio is a martingale under~$H_0$;
the cancellation that makes it work; non-negative martingale; drift under~$H_1$.
\end{keytakeaway}


% ══════════════════════════════════════════════════════════════════════════════
% ACT 5
% ══════════════════════════════════════════════════════════════════════════════
\newpage
\section{Act~5 --- Markov, Ville, and the Anytime-Valid Guarantee}

\subsection{The Simulation}

\begin{simulation}
This is the pivotal simulation of the entire presentation.

\textbf{Panel~1 --- Markov (single time):}
10,000 fair-coin paths of $\LR_n$.  A vertical line at $n=200$.
A threshold at $\LR = 1/\alpha = 20$ (for $\alpha = 0.05$).
Count: how many paths are above the threshold \emph{at step~200}?
Fraction $\leq 5\%$.

\textbf{Panel~2 --- Ville (any time):}
Same 10,000 paths.  Count: how many \emph{ever} crossed the threshold at
\emph{any} step?  Also $\leq 5\%$ --- but requires a stronger theorem.

\textbf{Panel~3 --- Comparison:}
10,000 paths of a standard $z$-score checked at every step.  Count how many
ever show significance.  Much higher than 5\% --- the peeking problem.

Summary:
\begin{itemize}[nosep]
  \item Markov: controls false positives at a single time \checkmark
  \item Ville: controls across all times simultaneously \checkmark
  \item Standard test with peeking: does NOT control $\times$
\end{itemize}
\end{simulation}

\subsection{Intuitive Explanation}

\begin{intuition}
We are about to prove the most important result in this entire presentation:
the \textbf{anytime-valid guarantee}.

Two steps:
\begin{enumerate}[nosep]
  \item \textbf{Markov's inequality} --- a simple bound at one point in time.
  \item \textbf{Ville's inequality} --- extends the bound to all points in
        time simultaneously.
\end{enumerate}
Ville's inequality is \emph{the reason} sequential testing works.
\end{intuition}

\subsection{Mathematical Formulation}

% ── Markov ─────────────────
\subsubsection*{Step 1: Markov's Inequality}

\begin{theorem}[Markov's Inequality]
For any non-negative random variable $X$ with finite $\E[X]$:
\[
  \boxed{\Prob(X \geq c) \;\leq\; \frac{\E[X]}{c}}
\]
\end{theorem}

\begin{plainlanguage}
``The probability that $X$ is at least $c$ cannot exceed the average of $X$
divided by $c$.''
\end{plainlanguage}

\textbf{Why is this true?  An intuitive proof.}

Suppose $X \geq 0$ always.  Let $q = \Prob(X \geq c)$ be the fraction of
outcomes at or above $c$.  Those outcomes contribute \emph{at least}
$q \times c$ to the average (each is $\geq c$).  The rest contribute $\geq 0$
(since $X \geq 0$):
\[
  \E[X] \;\geq\; q \cdot c + (1-q)\cdot 0 = q \cdot c
  \qquad\Longrightarrow\qquad
  q \leq \frac{\E[X]}{c}. \qquad\checkmark
\]

\textbf{Concrete example.}
Average income in a town is \texteuro50,000.
\[
  \Prob(\text{income} \geq \text{\texteuro}500{,}000)
  \;\leq\;
  \frac{50{,}000}{500{,}000} = 0.10 = 10\%.
\]
At most 10\% of people can earn \texteuro500,000+.

\textbf{Applying Markov to the likelihood ratio.}

Under $H_0$, $\LR_n$ is a martingale with $\E[\LR_n] = \LR_0 = 1$ and
$\LR_n \geq 0$.  Markov gives:
\[
  \Prob\!\left(\LR_n \geq \frac{1}{\alpha}\right)
  \leq \frac{\E[\LR_n]}{1/\alpha}
  = 1 \cdot \alpha
  = \alpha.
\]

\begin{plainlanguage}
At any single, pre-specified step $n$, the false positive probability
is at most $\alpha$.  Reassuring --- but only for \emph{one} moment.
\end{plainlanguage}

% ── Ville ──────────────────
\subsubsection*{Step 2: Ville's Inequality (1939)}

\begin{theorem}[Ville's Inequality]
For any \textbf{non-negative martingale} $\{M_n\}_{n \geq 0}$ with
$M_0 = m_0$:
\[
  \boxed{\Prob\!\left(\sup_{n \geq 0}\, M_n \;\geq\; c\right)
         \;\leq\; \frac{m_0}{c}}
\]
where $\sup_{n\geq 0} M_n$ is the highest value $M_n$ \emph{ever} reaches.
\end{theorem}

\textbf{Why this is remarkable:}

\begin{center}
\renewcommand{\arraystretch}{1.3}
\begin{tabular}{@{} l p{9cm} @{}}
\toprule
\textbf{Markov} & At any single \emph{fixed} moment, $\Prob(M_n \geq c)$ is small. \\
\textbf{Ville}  & $\Prob(M_n$ \textbf{ever} $\geq c$, at \textbf{any} moment$)$ is \emph{also} small. \\
\bottomrule
\end{tabular}
\end{center}

Markov would allow the process to spike above~$c$ at various times, as long as
it is rarely there at any single pre-specified time.  Ville says even those
spikes are collectively rare.

\textbf{The intuition behind Ville.}
For a non-negative martingale, reaching a high value ``uses up'' the available
expected value.  Once high, it must stay high on average (martingale property),
yet it cannot go below zero (non-negativity).  These two constraints together
limit how often it can ever be high.

\subsubsection*{Applying Ville to the likelihood ratio}

Set $M_n = \LR_n$, $m_0 = \LR_0 = 1$, $c = 1/\alpha$:
\[
  \boxed{\Prob\!\left(\LR_n \text{ ever reaches } \frac{1}{\alpha}
         \text{ or higher}\right) \;\leq\; \alpha}
\]

\begin{plainlanguage}
The probability that the likelihood ratio ever crosses our rejection
threshold --- at any point, across all flips, even if you run forever ---
is at most $\alpha$.

\textbf{This is the anytime-valid guarantee.}  You can peek after every
observation.  The false positive rate, across all checks combined, stays
$\leq \alpha$.
\end{plainlanguage}

\subsubsection*{Connecting back to Act~0}

Standard test statistics (like $p$-values from a $t$-test) are \emph{not}
martingales.  Checking repeatedly gives many independent chances to be fooled.

The likelihood ratio \emph{is} a martingale under~$H_0$.  Ville's inequality
converts this into a practical guarantee: \textbf{peeking is safe.}

\begin{keytakeaway}
\textbf{Key concepts:} Markov's inequality (with proof), Ville's inequality,
the anytime-valid guarantee, martingale + non-negative = peek-safe.
\end{keytakeaway}


% ══════════════════════════════════════════════════════════════════════════════
% ACT 6
% ══════════════════════════════════════════════════════════════════════════════
\newpage
\section{Act~6 --- SPRT: Wald's Sequential Probability Ratio Test (1945)}

\subsection{The Simulation}

\begin{simulation}
Two coins: one fair ($H_0$), one biased ($H_1$).  The user does not know
which they have.

$\LR_n$ is plotted as a path.  Two thresholds:
\begin{itemize}[nosep]
  \item \textbf{Upper} at $\LR_n = B$: cross this $\Rightarrow$ ``coin is biased.''
  \item \textbf{Lower} at $\LR_n = A$: fall below this $\Rightarrow$
        ``coin is fair.''
\end{itemize}
The path drifts toward one boundary and eventually crosses.

\textbf{User controls:} which coin is used, $\alpha$, $\beta$, $\delta$.
\end{simulation}

\subsection{Intuitive Explanation}

\begin{intuition}
Abraham Wald (1943, published 1945) asked: given that we can peek at the data
anytime, what is the most \emph{efficient} way to decide between two
hypotheses?

His answer: compute $\LR_n$ after each observation.  If it gets high enough,
conclude $H_1$.  Low enough, conclude $H_0$.  Otherwise, keep collecting.

Wald proved this is \textbf{optimal}: no other sequential test with the same
error guarantees needs fewer observations on average.
\end{intuition}

\subsection{Mathematical Formulation}

\subsubsection*{Decision rule}

Choose error tolerances $\alpha$ and $\beta$.
Compute two thresholds (\textbf{Wald's approximations} --- see note below):
\[
  \boxed{B = \frac{1-\beta}{\alpha}}, \qquad
  \boxed{A = \frac{\beta}{1-\alpha}}.
\]

\textbf{Example} ($\alpha = 0.05$, $\beta = 0.20$):
\[
  B = \frac{0.80}{0.05} = 16, \qquad
  A = \frac{0.20}{0.95} \approx 0.211.
\]

\begin{center}
\begin{tabular}{@{} l l l @{}}
\toprule
\textbf{Condition} & \textbf{Decision} & \textbf{Meaning} \\
\midrule
$\LR_n \geq B = 16$        & Reject $H_0$  & ``Coin is biased'' \\
$\LR_n \leq A \approx 0.21$  & Accept $H_0$  & ``Coin is fair'' \\
$A < \LR_n < B$            & Keep flipping & ``Not enough evidence yet'' \\
\bottomrule
\end{tabular}
\end{center}

\subsubsection*{A note on the thresholds}

The thresholds $B = (1-\beta)/\alpha$ and $A = \beta/(1-\alpha)$ are
\textbf{Wald's approximations}, not exact values.  Wald showed they are very
close to the true optimal boundaries and slightly conservative (actual error
rates are at or below nominal $\alpha$ and $\beta$).  Exact boundaries can be
computed numerically but rarely differ meaningfully.

\subsubsection*{Optimality of SPRT}

Wald \& Wolfowitz (1948) proved: among all sequential tests achieving error
rates $\leq \alpha$ and $\leq \beta$, the SPRT minimises the expected number
of observations under both $H_0$ and $H_1$.

\subsubsection*{The critical weakness}

SPRT requires specifying $\delta$ in advance.  If you set $\delta = 0.1$ but
the true bias is $0.03$, the likelihood ratio is computed against the wrong
alternative.  The test may take much longer, or fail to detect the effect.

In A/B testing --- where the effect size is typically unknown --- this is a
serious limitation.

\begin{keytakeaway}
\textbf{Key concepts:} SPRT decision rule, Wald's approximations, optimality,
critical dependence on a pre-specified effect size.
\end{keytakeaway}


% ══════════════════════════════════════════════════════════════════════════════
% ACT 7
% ══════════════════════════════════════════════════════════════════════════════
\newpage
\section{Act~7 --- mSPRT: The Mixture Sequential Probability Ratio Test}

\subsection{The Simulation}

\begin{simulation}
Like Act~6, but the user no longer specifies $\delta$.  A new slider
controls $\tau$ (width of a ``prior belief'' over effect sizes).

$\LR_n^{H}$ is plotted.  Only the \emph{upper} threshold $1/\alpha$ is
shown (no lower boundary).  Some runs cross early; some never cross.

\textbf{Comparison mode:} SPRT with wrong $\delta$ alongside mSPRT,
plus a \textbf{power curve} and \textbf{average stopping time chart}.
\end{simulation}

\subsection{Intuitive Explanation}

\begin{intuition}
SPRT's fatal flaw: you must know $\delta$ in advance.  The mSPRT's fix:
\textbf{bet on all effect sizes at once.}

Imagine a panel of experts, each with a different guess for $\delta$.  Each
runs their own SPRT.  You take a weighted average of their likelihood ratios.

Remarkably, this average is still a martingale under~$H_0$.  Ville's
inequality still applies.  You gain robustness without losing protection.
\end{intuition}

\subsection{Mathematical Formulation}

\subsubsection*{Why the average of martingales is a martingale}

If $\{A_n\}$ and $\{B_n\}$ are both martingales and $w \in [0,1]$ is a fixed
weight, let $C_n = w A_n + (1-w)B_n$.  Then:
\begin{align*}
  \E[C_n \given \text{past}]
  &= w\,\E[A_n \given \text{past}] + (1-w)\,\E[B_n \given \text{past}]
     & &\text{(linearity of $\E$)} \\
  &= w\,A_{n-1} + (1-w)\,B_{n-1}
     & &\text{(martingale property)} \\
  &= C_{n-1}. \qquad\checkmark
\end{align*}

\begin{plainlanguage}
The average of ``no trends'' is still ``no trend.''  Averaging does not break
the martingale property.
\end{plainlanguage}

\subsubsection*{The mixing idea}

Let $H$ be a \textbf{prior distribution} over effect sizes: a probability
distribution describing how plausible each $\delta$ is.  The mSPRT statistic
averages the SPRT likelihood ratio over all $\delta$ according to~$H$:
\[
  \boxed{\LR_n^{H} = \int \LR_n^{\delta} \, dH(\delta)}
\]

\begin{plainlanguage}
Symbol by symbol:
\begin{itemize}[nosep]
  \item $\LR_n^{\delta}$ = the likelihood ratio we would have computed
        assuming effect size $\delta$.
  \item $dH(\delta)$ = weighted by how probable $\delta$ is according to $H$.
  \item $\int$ = continuous version of summation --- averages over all
        possible $\delta$.
\end{itemize}
Think: ``average likelihood ratio across an infinitely large expert panel,
weighted by credibility.''
\end{plainlanguage}

\subsubsection*{Closed form for Normal observations}

Assumptions (used in practice by Eppo):
\begin{itemize}[nosep]
  \item Observations $x_i \sim \mathcal{N}(\mu, \sigma^2)$
        (Normal with mean $\mu$, known variance $\sigma^2$).
  \item Prior $H = \mathcal{N}(0, \tau^2)$ (Normal centred at~0,
        standard deviation $\tau$).
\end{itemize}

The integral has an exact algebraic solution:
\[
  \boxed{\LR_n^{H}
  = \sqrt{\frac{\sigma^2}{\sigma^2 + n\tau^2}}
    \;\exp\!\left(
      \frac{\tau^2\, n^2 \bar{x}_n^2}
           {2\,\sigma^2\,(\sigma^2 + n\tau^2)}
    \right)}
\]

\textbf{Unpacking every piece:}

\begin{center}
\renewcommand{\arraystretch}{1.4}
\begin{tabular}{@{} >{\raggedright}p{3cm} p{9cm} @{}}
\toprule
\textbf{Symbol} & \textbf{Meaning} \\
\midrule
$\sigma^2$ & Variance of individual observations (how noisy the data is).
  $\sigma$ is the standard deviation. \\
$n$ & Number of observations so far. \\
$\tau^2$ & Variance of the prior $H$.  Controls how wide a range of
  effect sizes we consider plausible. \\
$\bar{x}_n$ & Sample mean after $n$ observations:
  $\bar{x}_n = \frac{1}{n}\sum_{i=1}^n x_i$. \\
$n^2 \bar{x}_n^2$ & Equals $\left(\sum_{i=1}^n x_i\right)^2$: the square of
  the total sum.  Large when the cumulative signal is strong. \\
$\sigma^2 + n\tau^2$ & Combined scale.  As $n$ grows, $n\tau^2$ dominates. \\
$\sqrt{\frac{\sigma^2}{\sigma^2+n\tau^2}}$ & Shrinkage factor.  Slowly
  decreases as $n$ grows; prevents the statistic from growing merely because
  more data arrived. \\
$\exp(\cdot)$ & $e \approx 2.718$ raised to the argument.  Grows rapidly
  when $\bar{x}_n$ is far from zero. \\
\bottomrule
\end{tabular}
\end{center}

\begin{plainlanguage}
``The mSPRT score is large when the observed average is far from zero (strong
signal), when there are many observations, and when the noise is low relative
to the prior width.''
\end{plainlanguage}

\textbf{Practical note on $\sigma^2$:} The formula assumes $\sigma^2$ is
known.  In practice it is estimated from the data.  For $n \geq 100$ (almost
always true in A/B testing), estimation uncertainty is negligible.

\subsubsection*{Decision rule}

\[
  \text{Stop and reject } H_0 \text{ when:} \quad
  \boxed{\LR_n^{H} \;\geq\; \frac{1}{\alpha}}
\]

For $\alpha = 0.05$: stop when $\LR_n^{H} \geq 20$.

\textbf{No lower boundary.}
Unlike SPRT, mSPRT cannot definitively confirm $H_0$.  The mixture over all
effect sizes, including very small ones, means the statistic hovers near~1
indefinitely under $H_0$.  In practice, experiments run for a fixed maximum
duration; if the threshold is never crossed, conclude: ``No significant effect
detected.''

\subsubsection*{The role of $\tau$}

\begin{center}
\begin{tabular}{@{} l l l @{}}
\toprule
$\tau$ is\ldots & \textbf{Effect} & \textbf{Best for} \\
\midrule
Small ($\sim\!0.01$) & Conservative, slow to reject & Very small effects \\
Large ($\sim\!0.5$)  & Aggressive, fast for large effects & Large effects \\
Matched to true $\delta$ & Near-optimal & Good prior knowledge \\
\bottomrule
\end{tabular}
\end{center}

No universally optimal $\tau$ exists.  Calibrating $\tau$ from historical data
(GAVI, Act~8) gives a principled answer.

\subsubsection*{Ville's inequality still applies}

$\LR_n^{H}$ is a non-negative martingale under $H_0$ (mixture of
martingales), with $\LR_0^{H} = 1$.  Ville:
\[
  \boxed{\Prob\bigl(\LR_n^{H} \text{ ever reaches } 1/\alpha\bigr) \leq \alpha}
\]
The anytime-valid guarantee carries through unchanged.

\begin{keytakeaway}
\textbf{Key concepts:} mixture distribution, prior distribution, the role of
$\tau$, robustness to effect-size uncertainty, the closed-form formula,
practical notes ($\hat\sigma^2$, no lower boundary).
\end{keytakeaway}


% ══════════════════════════════════════════════════════════════════════════════
% ACT 8
% ══════════════════════════════════════════════════════════════════════════════
\newpage
\section{Act~8 --- Eppo's Implementation: CUPED, GAVI, \&\ Confidence
         Sequences}

This act has three independent sections, each addressing a different practical
concern: reducing noise (CUPED), calibrating the prior (GAVI), and reporting
results (confidence sequences).

\subsection{The Simulation}

\begin{simulation}
A full A/B test simulation.  Two groups (control / treatment), observations
arriving over time.

\textbf{Panel~1 --- CUPED toggle:}
mSPRT path with a toggle for variance reduction.  With CUPED on, the path is
less noisy and reaches the threshold sooner.

\textbf{Panel~2 --- Confidence sequence:}
A confidence band around the estimated treatment effect.  Starts wide, narrows
over time.  CUPED tightens it further.

\textbf{Panel~3 --- $\tau$ calibration:}
Average stopping time and power as functions of $\tau$.  A ``GAVI'' button
auto-calibrates from simulated historical experiments.

\textbf{User controls:} sample size, true effect size, $\tau$, correlation
$\rho$ (controls CUPED effectiveness).
\end{simulation}

% ── 8A: CUPED ──────────────
\subsection{Section 8A --- CUPED: Variance Reduction}

\subsubsection*{Intuitive Explanation}

\begin{intuition}
Testing whether a new feature increases purchases?  Pre-experiment data tells
you some users are heavy purchasers regardless.  That pre-existing variation is
noise.

CUPED (Controlled-experiment Using Pre-Experiment Data) removes this noise
before the test runs.  Like weighing yourself before and after a meal --- you
care about the change, not the baseline.

Less noise $\Rightarrow$ fewer observations needed.
\end{intuition}

\subsubsection*{Mathematical Formulation}

\textbf{Variables:}
\begin{itemize}[nosep]
  \item $Y$ = outcome during the experiment (e.g.\ purchase amount).
  \item $X$ = pre-experiment covariate (e.g.\ last month's purchases).
  \item $\E[X]$ = population average of $X$.
\end{itemize}

\textbf{CUPED-adjusted outcome:}
\[
  \boxed{Y^{*} = Y - \theta\,(X - \E[X])}
\]

\begin{plainlanguage}
\begin{itemize}[nosep]
  \item $\theta$: a coefficient chosen to remove as much noise as possible.
  \item $(X - \E[X])$: the centred pre-experiment variable (average = 0).
  \item $Y - \theta(X - \E[X])$: the outcome with the predictable component
        subtracted out.
\end{itemize}
\end{plainlanguage}

\textbf{Worked example.}
User purchased \texteuro80 during experiment ($Y=80$), \texteuro100 last
month ($X=100$), average is \texteuro70 ($\E[X]=70$), $\theta = 0.5$:
\[
  Y^* = 80 - 0.5 \times (100 - 70) = 80 - 15 = 65.
\]

\textbf{Optimal $\theta$:}
\[
  \boxed{\theta = \frac{\Cov(Y,X)}{\Var(X)}}
\]

\begin{plainlanguage}
This is the slope of a linear regression of $Y$ on $X$.
$\Cov(Y,X)$ measures how strongly $Y$ and $X$ move together.
$\Var(X)$ measures how spread out $X$ is.
\end{plainlanguage}

\textbf{Practical note:} $\theta$ is estimated from \emph{pre-experiment}
data, independent of treatment assignment.  This independence is why CUPED
does not compromise the test's validity.

\textbf{Variance reduction:}
\[
  \boxed{\Var(Y^*) = \Var(Y)\,(1 - \rho^2)}
\]
where $\rho = \Cov(Y,X) / \bigl(\SD(Y)\cdot\SD(X)\bigr)$ is the correlation
between $Y$ and~$X$.

\begin{center}
\begin{tabular}{@{} c c c r l @{}}
\toprule
$\rho$ & $\rho^2$ & $1-\rho^2$
       & \textbf{Variance remaining} & \textbf{Interpretation} \\
\midrule
0   & 0    & 1.00 & 100\% & No reduction \\
0.3 & 0.09 & 0.91 & 91\%  & Modest \\
0.5 & 0.25 & 0.75 & 75\%  & 25\% reduction \\
0.7 & 0.49 & 0.51 & 51\%  & Nearly half removed \\
0.9 & 0.81 & 0.19 & 19\%  & 81\% reduction --- dramatic \\
\bottomrule
\end{tabular}
\end{center}

\textbf{Derivation.}
$\Var(Y^*) = \Var(Y) + \theta^2\Var(X) - 2\theta\Cov(Y,X)$.
Substituting $\theta = \Cov(Y,X)/\Var(X)$ and simplifying yields
$\Var(Y)(1-\rho^2)$.

% ── 8B: GAVI ──────────────
\subsection{Section 8B --- GAVI: Calibrating $\tau$ from History}

\subsubsection*{Intuitive Explanation}

\begin{intuition}
How to choose $\tau$?  Look at past experiments.  If 100 previous A/B tests
showed effects in the range $\delta = 0.005$ to $0.10$, mostly around $0.02$,
set $\tau$ to match.

This is \textbf{GAVI} (Generalized Adjusted Variance Inflation): calibrate the
prior from empirical data. Near-optimal power for your platform's typical
effects.
\end{intuition}

\subsubsection*{Mathematical Formulation}

GAVI sets $H$ (the prior in the mSPRT integral) to match the empirical
distribution of past effect sizes:
\begin{enumerate}[nosep]
  \item Collect estimated effects $\delta_1, \ldots, \delta_k$ from $k$
        past experiments.
  \item Compute their standard deviation.
  \item Set $\tau$ equal (or proportional) to this standard deviation.
\end{enumerate}

% ── 8C: Confidence Sequences ──────────────
\subsection{Section 8C --- Confidence Sequences}

\subsubsection*{Intuitive Explanation}

\begin{intuition}
Standard tests report a \textbf{confidence interval}: a range of plausible
values for the effect (e.g.\ ``$+3\% \pm 2\%$'').

mSPRT gives something better: a \textbf{confidence sequence} --- a band that
updates after every observation and is valid at \emph{every} point in time.
It starts wide, narrows as data accumulates, and at any moment covers the true
effect with probability $\geq 1-\alpha$.
\end{intuition}

\subsubsection*{Mathematical Formulation}

The confidence sequence is obtained by \emph{inverting} the mSPRT rejection
rule.  At time $n$:
\[
  C(n) = \bigl\{\mu : \LR_n^{H}(\mu) < 1/\alpha \bigr\}
\]
\begin{plainlanguage}
$C(n)$ is the set of all parameter values $\mu$ for which the mSPRT,
computed as if $\mu$ were the null, does not exceed the threshold.
\end{plainlanguage}

For the Normal case, this yields an interval:
\[
  \bar{x}_n \;\pm\; \text{margin}(n, \sigma^2, \tau^2, \alpha)
\]
The margin shrinks as $n$ grows.

The guarantee:
\[
  \boxed{\Prob\!\bigl(\text{true } \mu \in C(n) \text{ for all } n
         \text{ simultaneously}\bigr) \geq 1-\alpha}
\]

\begin{plainlanguage}
The band covers the true parameter at all times --- not just at one
pre-specified moment.  This is the confidence-sequence counterpart of Ville's
inequality.
\end{plainlanguage}

% ── 8D: Full Pipeline ──────────────
\subsection{Section 8D --- The Full Pipeline}

\begin{enumerate}
  \item \textbf{Before the experiment:} Collect pre-experiment covariate $X$.
        Estimate $\theta$ from historical data.
  \item \textbf{During:} For each observation $Y$, compute
        $Y^* = Y - \theta(X - \E[X])$.
  \item \textbf{Calibrate $\tau$} from historical effect sizes (GAVI).
  \item \textbf{Run mSPRT} on $Y^*$ with parameter $\tau$, threshold
        $1/\alpha$.
  \item \textbf{Monitor} the confidence sequence.
  \item \textbf{Decide:} If $\LR_n^{H} \geq 1/\alpha$, reject~$H_0$.
        Otherwise, at maximum planned sample size, conclude ``no significant
        effect detected.''
\end{enumerate}

\subsubsection*{Why CUPED does not break the mSPRT guarantee}

The mSPRT is applied to $Y^*$ instead of~$Y$.  Same formulas, smaller
$\sigma^2$.  Crucially:
\begin{enumerate}[nosep]
  \item $\theta$ is estimated from \emph{pre-experiment} data (independent
        of treatment).
  \item CUPED is applied \emph{before} the mSPRT statistic is computed.
  \item The martingale property depends only on correct modelling of the
        observations, which holds for the reduced-variance $Y^*$.
\end{enumerate}
Ville's inequality still holds.  The test is still anytime-valid.

\begin{keytakeaway}
\textbf{Key concepts:} CUPED, pre-experiment covariate, variance reduction,
correlation, GAVI, $\tau$ calibration, confidence sequences, the full Eppo
pipeline.

The test is \textbf{anytime-valid}, \textbf{variance-reduced}, and
\textbf{robust to effect-size uncertainty}.
\end{keytakeaway}


% ══════════════════════════════════════════════════════════════════════════════
\newpage
\section{Navigation and UX Design}

\subsection{Between Acts --- Transition Cards}

Each transition features a brief interstitial card with a bridge sentence and
a ``Key insight so far'' box.

\begin{center}
\renewcommand{\arraystretch}{1.3}
\begin{tabular}{@{} l p{9cm} @{}}
\toprule
\textbf{Transition} & \textbf{Bridge sentence} \\
\midrule
Act~0 $\to$ 1 & ``To build a peek-safe test, we first need the right
  building block.  It starts with a coin flip.'' \\
Act~1 $\to$ 2 & ``Now let's put money on each step --- and learn why no
  strategy can beat a fair game.'' \\
Act~2 $\to$ 3 & ``Fair games can't be beaten.  Now let's build a detective
  --- a tool for telling two coins apart.'' \\
Act~3 $\to$ 4 & ``We can measure evidence with a likelihood ratio.  But
  here's the surprise: that evidence score is itself a martingale.'' \\
Act~4 $\to$ 5 & ``The likelihood ratio is a martingale.  Now let's see why
  that single fact gives us the peeking guarantee.'' \\
Act~5 $\to$ 6 & ``We have the theory: peeking is safe.  Now let's turn it
  into a practical decision procedure.'' \\
Act~6 $\to$ 7 & ``SPRT is optimal but inflexible.  What if we don't know the
  effect size?'' \\
Act~7 $\to$ 8 & ``mSPRT gives a robust, peek-safe test.  Now let's add
  Eppo's practical enhancements.'' \\
\bottomrule
\end{tabular}
\end{center}

\subsection{Persistent Controls}

A slider panel available throughout all acts:
\begin{itemize}[nosep]
  \item Probability $p$ (or effect size $\delta$ from Act~3 onward)
  \item Significance level $\alpha$
  \item Number of steps / observations $n$
  \item Simulation speed
  \item Number of simultaneous runs
\end{itemize}

\subsection{Persistent Sidebar}

The notation cheat sheet (Section~\ref{sec:notation}) should be accessible as a
collapsible sidebar from every act.

\subsection{Visual Design Principles}

\begin{itemize}[nosep]
  \item The same visual object (path / process) carries through all eight
        acts, relabelled at each stage.
  \item Threshold lines appear gradually: none in Acts~0--1; $\pm\sqrt{n}$
        envelope in Act~1; both SPRT thresholds in Act~6; only the upper
        threshold from Act~7; confidence bands in Act~8.
  \item Formula panels slide in/out; not shown by default on first reading.
  \item Every formula has a plain-language translation directly below it.
  \item Key terms on first use are bolded and defined in-line.
\end{itemize}


% ══════════════════════════════════════════════════════════════════════════════
\newpage
\section{Summary: The Through-Line}

\begin{center}
\renewcommand{\arraystretch}{1.3}
\begin{tabular}{@{} c l l l @{}}
\toprule
\textbf{Act} & \textbf{Object} & \textbf{Represents} & \textbf{Key idea} \\
\midrule
0 & Repeated $p$-values
  & The peeking problem
  & Checking repeatedly inflates FP \\
1 & Random walk $S_n$
  & Position
  & $\SD$ grows as $\sqrt{n}$ \\
2 & Winnings $M_n$
  & Gambler's fortune
  & Martingale; Doob: $\E[M_\tau]=M_0$ \\
3 & Likelihood ratio $\LR_n$
  & Evidence against $H_0$
  & Comparing hypotheses flip by flip \\
4 & $\LR_n$ under $H_0$
  & LR is a martingale
  & The cancellation proof; $\LR_0{=}1$, $\LR_n{\geq}0$ \\
5 & Markov $+$ Ville
  & Anytime-valid guarantee
  & $\Prob(\LR_n \text{ ever}{\geq}1/\alpha)\leq\alpha$ \\
6 & SPRT
  & Wald's decision rule
  & Optimal but requires known $\delta$ \\
7 & $\LR_n^{H}$
  & Robust evidence
  & Mixture of martingales is a martingale \\
8 & CUPED + CS
  & Less noisy, fully reported
  & Same guarantee, faster detection \\
\bottomrule
\end{tabular}
\end{center}

\bigskip
\textbf{The core message:} All eight acts build on the same mathematical
object --- a non-negative martingale --- viewed through increasingly powerful
lenses.  The anytime-valid guarantee flows from a single elegant fact:
\[
  \textit{A fair game cannot be beaten by choosing when to quit.}
\]


% ══════════════════════════════════════════════════════════════════════════════
\newpage
\section{References}
\label{sec:references}

\begin{enumerate}
  \item \textbf{Ville, J.} (1939).
        \textit{\'Etude critique de la notion de collectif}.
        Gauthier-Villars, Paris.
        --- Original proof of Ville's inequality for non-negative
        supermartingales.

  \item \textbf{Wald, A.} (1945).
        Sequential tests of statistical hypotheses.
        \textit{Annals of Mathematical Statistics}, 16(2), 117--186.
        \url{https://doi.org/10.1214/aoms/1177731118}
        --- Foundational paper on the SPRT.

  \item \textbf{Wald, A.\ \& Wolfowitz, J.} (1948).
        Optimum character of the sequential probability ratio test.
        \textit{Annals of Mathematical Statistics}, 19(3), 326--339.
        \url{https://doi.org/10.1214/aoms/1177730197}
        --- Proof that SPRT minimises expected sample size.

  \item \textbf{Doob, J.\,L.} (1953).
        \textit{Stochastic Processes}.
        Wiley, New York.
        --- Definitive treatment of martingale theory including the
        Optional Stopping Theorem.

  \item \textbf{Robbins, H.} (1970).
        Statistical methods related to the law of the iterated logarithm.
        \textit{Annals of Mathematical Statistics}, 41(5), 1397--1409.
        \url{https://doi.org/10.1214/aoms/1177696786}
        --- Early work on mixture sequential tests.

  \item \textbf{Deng, A., Xu, Y., Kohavi, R., \& Walker, T.} (2013).
        Improving the sensitivity of online controlled experiments by
        utilizing pre-experiment data.
        \textit{Proc.\ WSDM}, 123--132.
        \url{https://doi.org/10.1145/2433396.2433413}
        --- The original CUPED paper.

  \item \textbf{Johari, R., Pekelis, L., \& Walsh, D.J.} (2017).
        Always valid inference: Continuous monitoring of A/B tests.
        \textit{Operations Research}, 65(6), 1651--1667.
        \url{https://doi.org/10.1287/opre.2017.1607}
        --- mSPRT for A/B testing.

  \item \textbf{Howard, S.\,R., Ramdas, A., McAuliffe, J., \& Sekhon, J.}
        (2021).
        Time-uniform, nonparametric, nonasymptotic confidence sequences.
        \textit{Annals of Statistics}, 49(2), 1055--1080.
        \url{https://doi.org/10.1214/20-AOS1991}
        --- Theory of confidence sequences and anytime-valid inference.

  \item \textbf{Schmit, M.\ \& Miller, T.} (2024).
        Sequential confidence intervals for comparing two proportions with
        variance reduction.
        Eppo (available in \texttt{literature/} folder).
        --- Eppo's implementation: CUPED-based variance reduction, confidence
        sequences, and GAVI.
\end{enumerate}


% ══════════════════════════════════════════════════════════════════════════════
\newpage
\section*{Appendix: What Changed from v2}
\addcontentsline{toc}{section}{Appendix: What Changed from v2}

\begin{enumerate}[nosep]
  \item Converted from Markdown to \LaTeX{} for proper mathematical
        typesetting: fractions, summation limits, aligned derivations,
        boxed key formulas, cancel notation in the martingale proof, etc.
  \item All formulas now use standard \LaTeX{} math mode instead of Unicode
        approximations.
  \item Added coloured boxes to visually distinguish simulations, intuitive
        explanations, and key takeaways.
  \item Added a table of contents for easy navigation.
  \item References now include DOI links where available.
  \item No content changes --- the pedagogical structure and all text are
        preserved from v2.
\end{enumerate}

\end{document}
